% Options for packages loaded elsewhere
\PassOptionsToPackage{unicode}{hyperref}
\PassOptionsToPackage{hyphens}{url}
\documentclass[
]{article}
\usepackage{xcolor}
\usepackage[margin=1in]{geometry}
\usepackage{amsmath,amssymb}
\setcounter{secnumdepth}{-\maxdimen} % remove section numbering
\usepackage{iftex}
\ifPDFTeX
  \usepackage[T1]{fontenc}
  \usepackage[utf8]{inputenc}
  \usepackage{textcomp} % provide euro and other symbols
\else % if luatex or xetex
  \usepackage{unicode-math} % this also loads fontspec
  \defaultfontfeatures{Scale=MatchLowercase}
  \defaultfontfeatures[\rmfamily]{Ligatures=TeX,Scale=1}
\fi
\usepackage{lmodern}
\ifPDFTeX\else
  % xetex/luatex font selection
\fi
% Use upquote if available, for straight quotes in verbatim environments
\IfFileExists{upquote.sty}{\usepackage{upquote}}{}
\IfFileExists{microtype.sty}{% use microtype if available
  \usepackage[]{microtype}
  \UseMicrotypeSet[protrusion]{basicmath} % disable protrusion for tt fonts
}{}
\makeatletter
\@ifundefined{KOMAClassName}{% if non-KOMA class
  \IfFileExists{parskip.sty}{%
    \usepackage{parskip}
  }{% else
    \setlength{\parindent}{0pt}
    \setlength{\parskip}{6pt plus 2pt minus 1pt}}
}{% if KOMA class
  \KOMAoptions{parskip=half}}
\makeatother
\usepackage{color}
\usepackage{fancyvrb}
\newcommand{\VerbBar}{|}
\newcommand{\VERB}{\Verb[commandchars=\\\{\}]}
\DefineVerbatimEnvironment{Highlighting}{Verbatim}{commandchars=\\\{\}}
% Add ',fontsize=\small' for more characters per line
\usepackage{framed}
\definecolor{shadecolor}{RGB}{248,248,248}
\newenvironment{Shaded}{\begin{snugshade}}{\end{snugshade}}
\newcommand{\AlertTok}[1]{\textcolor[rgb]{0.94,0.16,0.16}{#1}}
\newcommand{\AnnotationTok}[1]{\textcolor[rgb]{0.56,0.35,0.01}{\textbf{\textit{#1}}}}
\newcommand{\AttributeTok}[1]{\textcolor[rgb]{0.13,0.29,0.53}{#1}}
\newcommand{\BaseNTok}[1]{\textcolor[rgb]{0.00,0.00,0.81}{#1}}
\newcommand{\BuiltInTok}[1]{#1}
\newcommand{\CharTok}[1]{\textcolor[rgb]{0.31,0.60,0.02}{#1}}
\newcommand{\CommentTok}[1]{\textcolor[rgb]{0.56,0.35,0.01}{\textit{#1}}}
\newcommand{\CommentVarTok}[1]{\textcolor[rgb]{0.56,0.35,0.01}{\textbf{\textit{#1}}}}
\newcommand{\ConstantTok}[1]{\textcolor[rgb]{0.56,0.35,0.01}{#1}}
\newcommand{\ControlFlowTok}[1]{\textcolor[rgb]{0.13,0.29,0.53}{\textbf{#1}}}
\newcommand{\DataTypeTok}[1]{\textcolor[rgb]{0.13,0.29,0.53}{#1}}
\newcommand{\DecValTok}[1]{\textcolor[rgb]{0.00,0.00,0.81}{#1}}
\newcommand{\DocumentationTok}[1]{\textcolor[rgb]{0.56,0.35,0.01}{\textbf{\textit{#1}}}}
\newcommand{\ErrorTok}[1]{\textcolor[rgb]{0.64,0.00,0.00}{\textbf{#1}}}
\newcommand{\ExtensionTok}[1]{#1}
\newcommand{\FloatTok}[1]{\textcolor[rgb]{0.00,0.00,0.81}{#1}}
\newcommand{\FunctionTok}[1]{\textcolor[rgb]{0.13,0.29,0.53}{\textbf{#1}}}
\newcommand{\ImportTok}[1]{#1}
\newcommand{\InformationTok}[1]{\textcolor[rgb]{0.56,0.35,0.01}{\textbf{\textit{#1}}}}
\newcommand{\KeywordTok}[1]{\textcolor[rgb]{0.13,0.29,0.53}{\textbf{#1}}}
\newcommand{\NormalTok}[1]{#1}
\newcommand{\OperatorTok}[1]{\textcolor[rgb]{0.81,0.36,0.00}{\textbf{#1}}}
\newcommand{\OtherTok}[1]{\textcolor[rgb]{0.56,0.35,0.01}{#1}}
\newcommand{\PreprocessorTok}[1]{\textcolor[rgb]{0.56,0.35,0.01}{\textit{#1}}}
\newcommand{\RegionMarkerTok}[1]{#1}
\newcommand{\SpecialCharTok}[1]{\textcolor[rgb]{0.81,0.36,0.00}{\textbf{#1}}}
\newcommand{\SpecialStringTok}[1]{\textcolor[rgb]{0.31,0.60,0.02}{#1}}
\newcommand{\StringTok}[1]{\textcolor[rgb]{0.31,0.60,0.02}{#1}}
\newcommand{\VariableTok}[1]{\textcolor[rgb]{0.00,0.00,0.00}{#1}}
\newcommand{\VerbatimStringTok}[1]{\textcolor[rgb]{0.31,0.60,0.02}{#1}}
\newcommand{\WarningTok}[1]{\textcolor[rgb]{0.56,0.35,0.01}{\textbf{\textit{#1}}}}
\usepackage{graphicx}
\makeatletter
\newsavebox\pandoc@box
\newcommand*\pandocbounded[1]{% scales image to fit in text height/width
  \sbox\pandoc@box{#1}%
  \Gscale@div\@tempa{\textheight}{\dimexpr\ht\pandoc@box+\dp\pandoc@box\relax}%
  \Gscale@div\@tempb{\linewidth}{\wd\pandoc@box}%
  \ifdim\@tempb\p@<\@tempa\p@\let\@tempa\@tempb\fi% select the smaller of both
  \ifdim\@tempa\p@<\p@\scalebox{\@tempa}{\usebox\pandoc@box}%
  \else\usebox{\pandoc@box}%
  \fi%
}
% Set default figure placement to htbp
\def\fps@figure{htbp}
\makeatother
\setlength{\emergencystretch}{3em} % prevent overfull lines
\providecommand{\tightlist}{%
  \setlength{\itemsep}{0pt}\setlength{\parskip}{0pt}}
\usepackage{bookmark}
\IfFileExists{xurl.sty}{\usepackage{xurl}}{} % add URL line breaks if available
\urlstyle{same}
\hypersetup{
  pdftitle={DEshort},
  hidelinks,
  pdfcreator={LaTeX via pandoc}}

\title{DEshort}
\author{}
\date{\vspace{-2.5em}}

\begin{document}
\maketitle

Differential expression analysis aims at distinguish the cluster of
genes that are over- or under- represented in each sample. Comparing
case and control samples aims at identifying which are the genes
correlated to the disease.

\subsubsection{0. importing data}\label{importing-data}

Import the dataframe and filter unuseful information

\begin{Shaded}
\begin{Highlighting}[]
\NormalTok{data }\OtherTok{\textless{}{-}} \FunctionTok{read.table}\NormalTok{(}\StringTok{"R\_steps/reads\_gene\_counts.txt"}\NormalTok{, }
                   \AttributeTok{header =} \ConstantTok{TRUE}\NormalTok{, }
                   \AttributeTok{sep =} \StringTok{"}\SpecialCharTok{\textbackslash{}t}\StringTok{"}\NormalTok{)}
\FunctionTok{colnames}\NormalTok{(data)}
\end{Highlighting}
\end{Shaded}

\begin{verbatim}
##  [1] "Geneid"     "Chr"        "Start"      "End"        "Strand"    
##  [6] "Length"     "SRR7821918" "SRR7821919" "SRR7821920" "SRR7821921"
## [11] "SRR7821922" "SRR7821923" "SRR7821924" "SRR7821925" "SRR7821927"
## [16] "SRR7821937" "SRR7821938" "SRR7821939" "SRR7821940" "SRR7821941"
## [21] "SRR7821942"
\end{verbatim}

\begin{Shaded}
\begin{Highlighting}[]
\FunctionTok{dim}\NormalTok{(data)}
\end{Highlighting}
\end{Shaded}

\begin{verbatim}
## [1] 78334    21
\end{verbatim}

results in a table with 78334 obs for 21 variables (15 of which are our
samples). Then, filter for the columns corresponding to the samples and
ignore the rest of the information. Specifically, remove the columns
containing Chr, Start, End, Strand and Length.

\begin{Shaded}
\begin{Highlighting}[]
\NormalTok{data }\OtherTok{\textless{}{-}}\NormalTok{ data[,}\SpecialCharTok{{-}}\DecValTok{2}\SpecialCharTok{:{-}}\DecValTok{6}\NormalTok{]}
\FunctionTok{colnames}\NormalTok{(data)}
\end{Highlighting}
\end{Shaded}

\begin{verbatim}
##  [1] "Geneid"     "SRR7821918" "SRR7821919" "SRR7821920" "SRR7821921"
##  [6] "SRR7821922" "SRR7821923" "SRR7821924" "SRR7821925" "SRR7821927"
## [11] "SRR7821937" "SRR7821938" "SRR7821939" "SRR7821940" "SRR7821941"
## [16] "SRR7821942"
\end{verbatim}

Create a matrix ready to feed into DESeq2. It has to be ``numeric''.

\begin{Shaded}
\begin{Highlighting}[]
\FunctionTok{rownames}\NormalTok{(data) }\OtherTok{\textless{}{-}}\NormalTok{ data[,}\DecValTok{1}\NormalTok{]   }\CommentTok{\# set first column as row names}
\NormalTok{data }\OtherTok{\textless{}{-}}\NormalTok{ data[ , }\SpecialCharTok{{-}}\DecValTok{1}\NormalTok{]          }\CommentTok{\# remove the first column}
\end{Highlighting}
\end{Shaded}

\subsection{5. Exploratory data analysis
(DESeq2)}\label{exploratory-data-analysis-deseq2}

Install DESeq2

\begin{verbatim}
if (!require("BiocManager", quietly = TRUE))
    install.packages("BiocManager")

BiocManager::install("DESeq2")
\end{verbatim}

Get the sample information into a table containing the sample name and
to which group it belongs.

\begin{Shaded}
\begin{Highlighting}[]
\NormalTok{sample\_info }\OtherTok{\textless{}{-}} \FunctionTok{read.delim}\NormalTok{(}
  \StringTok{"R\_steps/sample\_info.txt"}\NormalTok{,}
  \AttributeTok{header =} \ConstantTok{TRUE}\NormalTok{,}
  \AttributeTok{sep =} \StringTok{" "}\NormalTok{,}
  \AttributeTok{stringsAsFactors =} \ConstantTok{FALSE}\NormalTok{,}
  \AttributeTok{check.names =} \ConstantTok{FALSE}
\NormalTok{)}
\end{Highlighting}
\end{Shaded}

Now give the data to DEseq2 so that it can create a dds (DESeq Data Set
that is created from the data matrix).

\begin{Shaded}
\begin{Highlighting}[]
\NormalTok{dds }\OtherTok{\textless{}{-}}\NormalTok{ DESeq2}\SpecialCharTok{::}\FunctionTok{DESeqDataSetFromMatrix}\NormalTok{(}\AttributeTok{countData =}\NormalTok{ data,}
                               \AttributeTok{colData =}\NormalTok{ sample\_info,}
                               \AttributeTok{design=} \SpecialCharTok{\textasciitilde{}}\NormalTok{ Group)}
\end{Highlighting}
\end{Shaded}

Now remove the dependence of the variance on the mean by running
\texttt{DESeq2::vst()} on dds. For visualization, we typically do not
want to consider the experimental groups so set \texttt{blind\ =\ TRUE}.

plot the PCA to check whether it is all right with the data and it is
possible to distinguish in the 4 clusters where they belong to.

\begin{Shaded}
\begin{Highlighting}[]
\CommentTok{\# PCA using top 500 features by variance}
\NormalTok{DESeq2}\SpecialCharTok{::}\FunctionTok{plotPCA}\NormalTok{(vsd, }\AttributeTok{intgroup =} \StringTok{"Group"}\NormalTok{)}
\end{Highlighting}
\end{Shaded}

\begin{verbatim}
## using ntop=500 top features by variance
\end{verbatim}

\pandocbounded{\includegraphics[keepaspectratio]{DEshort_files/figure-latex/unnamed-chunk-7-1.pdf}}

\subsection{6. Differential Expression
analysis}\label{differential-expression-analysis}

\begin{Shaded}
\begin{Highlighting}[]
\NormalTok{dss }\OtherTok{\textless{}{-}}\NormalTok{ DESeq2}\SpecialCharTok{::}\FunctionTok{DESeq}\NormalTok{(dds)}
\end{Highlighting}
\end{Shaded}

\begin{verbatim}
## estimating size factors
\end{verbatim}

\begin{verbatim}
## estimating dispersions
\end{verbatim}

\begin{verbatim}
## gene-wise dispersion estimates
\end{verbatim}

\begin{verbatim}
## mean-dispersion relationship
\end{verbatim}

\begin{verbatim}
## final dispersion estimates
\end{verbatim}

\begin{verbatim}
## fitting model and testing
\end{verbatim}

\begin{Shaded}
\begin{Highlighting}[]
\NormalTok{dss}
\end{Highlighting}
\end{Shaded}

\begin{verbatim}
## class: DESeqDataSet 
## dim: 78334 15 
## metadata(1): version
## assays(4): counts mu H cooks
## rownames(78334): ENSMUSG00000104478 ENSMUSG00000104385 ...
##   ENSMUSG00002074899 ENSMUSG00002076890
## rowData names(30): baseMean baseVar ... deviance maxCooks
## colnames(15): SRR7821918 SRR7821919 ... SRR7821941 SRR7821942
## colData names(3): Sample Group sizeFactor
\end{verbatim}

\begin{Shaded}
\begin{Highlighting}[]
\NormalTok{DESeq2}\SpecialCharTok{::}\FunctionTok{resultsNames}\NormalTok{(dss)}
\end{Highlighting}
\end{Shaded}

\begin{verbatim}
## [1] "Intercept"                              
## [2] "Group_Lung_DKO_Control_vs_Lung_DKO_Case"
## [3] "Group_Lung_WT_Case_vs_Lung_DKO_Case"    
## [4] "Group_Lung_WT_Control_vs_Lung_DKO_Case"
\end{verbatim}

\subsubsection{check resultsNames and extract the comparisons of
interest}\label{check-resultsnames-and-extract-the-comparisons-of-interest}

To count how many are differentially expressed, check how many are not
NA in the padj column and lower than the classic 0.05 threshold.

Check the number of differentially expressed genes in each comparison
with \emph{padj \textless{} 0.05}

\begin{Shaded}
\begin{Highlighting}[]
\NormalTok{get\_results }\OtherTok{\textless{}{-}} \ControlFlowTok{function}\NormalTok{(DEout, this, to\_that)\{}
\NormalTok{  res\_comp }\OtherTok{\textless{}{-}}\NormalTok{ DESeq2}\SpecialCharTok{::}\FunctionTok{results}\NormalTok{(DEout, }\AttributeTok{contrast =} \FunctionTok{c}\NormalTok{(}\StringTok{"Group"}\NormalTok{, this, to\_that)) }\CommentTok{\# extract comparison}
  
\NormalTok{  filt\_comp }\OtherTok{\textless{}{-}}\NormalTok{ res\_comp[res\_comp}\SpecialCharTok{$}\NormalTok{padj }\SpecialCharTok{\textless{}} \FloatTok{0.05} \SpecialCharTok{\&} \SpecialCharTok{!} \FunctionTok{is.na}\NormalTok{(res\_comp}\SpecialCharTok{$}\NormalTok{padj), ] }\CommentTok{\# filter for padj and remove NA}
  
  \FunctionTok{return}\NormalTok{(filt\_comp)}
\NormalTok{\}}

\NormalTok{WT\_Case\_DKO\_Case }\OtherTok{\textless{}{-}} \FunctionTok{get\_results}\NormalTok{(dss, }\StringTok{"Lung\_WT\_Case"}\NormalTok{, }\StringTok{"Lung\_DKO\_Case"}\NormalTok{) }\CommentTok{\# WT cases against DKO cases}
\NormalTok{WT\_Control\_WT\_Case }\OtherTok{\textless{}{-}} \FunctionTok{get\_results}\NormalTok{(dss, }\StringTok{"Lung\_WT\_Control"}\NormalTok{, }\StringTok{"Lung\_WT\_Case"}\NormalTok{) }\CommentTok{\# WT control against WT case}
\NormalTok{DKO\_Control\_DKO\_Case }\OtherTok{\textless{}{-}} \FunctionTok{get\_results}\NormalTok{(dss, }\StringTok{"Lung\_DKO\_Control"}\NormalTok{, }\StringTok{"Lung\_DKO\_Case"}\NormalTok{) }\CommentTok{\# DKO control against DKO case}
\end{Highlighting}
\end{Shaded}

\subsubsection{Heatmap}\label{heatmap}

Produce heatmaps to visualize the differences in the expression between
groups.

\begin{Shaded}
\begin{Highlighting}[]
\CommentTok{\# 1. Extract VSD matrix}
\NormalTok{vsd\_mat }\OtherTok{\textless{}{-}}\NormalTok{ SummarizedExperiment}\SpecialCharTok{::}\FunctionTok{assay}\NormalTok{(vsd)}

\NormalTok{heatplot }\OtherTok{\textless{}{-}} \ControlFlowTok{function}\NormalTok{(comparison, samples, vsd, title, }\AttributeTok{leg =} \ConstantTok{TRUE}\NormalTok{)\{}
  \CommentTok{\# 2. Select significant genes for the comparison}
\NormalTok{  filtered }\OtherTok{\textless{}{-}}\NormalTok{ comparison[(comparison}\SpecialCharTok{$}\NormalTok{log2FoldChange }\SpecialCharTok{\textgreater{}} \DecValTok{1} \SpecialCharTok{|}\NormalTok{ comparison}\SpecialCharTok{$}\NormalTok{log2FoldChange }\SpecialCharTok{\textless{}} \SpecialCharTok{{-}}\DecValTok{1}\NormalTok{), ] }\CommentTok{\# | log2FC | \textgreater{} 1 filtering}
\NormalTok{  filtered }\OtherTok{\textless{}{-}}\NormalTok{ filtered[}\FunctionTok{order}\NormalTok{(filtered}\SpecialCharTok{$}\NormalTok{padj),] }\CommentTok{\# increasing order adjusted p{-}val}
\NormalTok{  filtered }\OtherTok{\textless{}{-}}\NormalTok{ filtered[}\DecValTok{1}\SpecialCharTok{:}\DecValTok{25}\NormalTok{,] }\CommentTok{\# select the first 25 entries}
\NormalTok{  heat\_matrix }\OtherTok{\textless{}{-}}\NormalTok{ vsd[}\FunctionTok{rownames}\NormalTok{(filtered),] }\CommentTok{\# obtain a DESeq object matrix by filtering the VST result object }
  
  \CommentTok{\# 3. Build annotation data frame for the comparison}
\NormalTok{  ann }\OtherTok{\textless{}{-}}\NormalTok{ samples[, }\StringTok{"Group"}\NormalTok{, drop }\OtherTok{=} \ConstantTok{FALSE}\NormalTok{] }\CommentTok{\# extract Group column for the samples information}
  \FunctionTok{rownames}\NormalTok{(ann) }\OtherTok{\textless{}{-}}\NormalTok{ samples}\SpecialCharTok{$}\NormalTok{Sample }\CommentTok{\# index ann based on the samples}
  
  \CommentTok{\# 4. Plot heatmap}
  \FunctionTok{library}\NormalTok{(pheatmap)}
  \FunctionTok{pheatmap}\NormalTok{(heat\_matrix,}
           \AttributeTok{annotation\_col =}\NormalTok{ ann,}
           \AttributeTok{show\_rownames =} \ConstantTok{TRUE}\NormalTok{,}
           \AttributeTok{fontsize\_col =} \DecValTok{10}\NormalTok{,}
           \AttributeTok{clustering\_distance\_rows =} \StringTok{"euclidean"}\NormalTok{,}
           \AttributeTok{clustering\_distance\_cols =} \StringTok{"euclidean"}\NormalTok{,}
           \AttributeTok{clustering\_method =} \StringTok{"complete"}\NormalTok{,}
           \AttributeTok{scale =} \StringTok{"row"}\NormalTok{,}
           \AttributeTok{main =}\NormalTok{ title,}
           \AttributeTok{annotation\_legend =}\NormalTok{ leg,}
           \AttributeTok{show\_colnames =} \ConstantTok{FALSE}\NormalTok{)}
\NormalTok{\}}

\FunctionTok{heatplot}\NormalTok{(WT\_Case\_DKO\_Case, sample\_info, vsd\_mat, }\AttributeTok{title =} \StringTok{"WT vs DKO cases"}\NormalTok{)}
\end{Highlighting}
\end{Shaded}

\pandocbounded{\includegraphics[keepaspectratio]{DEshort_files/figure-latex/unnamed-chunk-10-1.pdf}}

\begin{Shaded}
\begin{Highlighting}[]
\FunctionTok{heatplot}\NormalTok{(WT\_Control\_WT\_Case, sample\_info, vsd\_mat, }\AttributeTok{title =} \StringTok{"WT controls vs cases"}\NormalTok{)}
\end{Highlighting}
\end{Shaded}

\pandocbounded{\includegraphics[keepaspectratio]{DEshort_files/figure-latex/unnamed-chunk-10-2.pdf}}

\begin{Shaded}
\begin{Highlighting}[]
\FunctionTok{heatplot}\NormalTok{(DKO\_Control\_DKO\_Case, sample\_info, vsd\_mat, }\AttributeTok{title =} \StringTok{"DKO controls vs cases"}\NormalTok{)}
\end{Highlighting}
\end{Shaded}

\pandocbounded{\includegraphics[keepaspectratio]{DEshort_files/figure-latex/unnamed-chunk-10-3.pdf}}

\textbf{SUMMARY REPORT TABLE 1}

How many genes are differentially expressed (DE) in the pairwise
comparisons you selected (e.g.~padj \textless{} 0.05)?

\begin{Shaded}
\begin{Highlighting}[]
\FunctionTok{paste0}\NormalTok{(}\StringTok{"There has been found "}\NormalTok{, }\FunctionTok{nrow}\NormalTok{(WT\_Case\_DKO\_Case), }\StringTok{" DE genes in WT vs DKO case comparison having padj \textless{} 0.05."}\NormalTok{)}
\end{Highlighting}
\end{Shaded}

\begin{verbatim}
## [1] "There has been found 7976 DE genes in WT vs DKO case comparison having padj < 0.05."
\end{verbatim}

\begin{Shaded}
\begin{Highlighting}[]
\FunctionTok{paste0}\NormalTok{(}\StringTok{"There has been found "}\NormalTok{, }\FunctionTok{nrow}\NormalTok{(WT\_Case\_DKO\_Case[}\FunctionTok{which}\NormalTok{(WT\_Case\_DKO\_Case}\SpecialCharTok{$}\NormalTok{log2FoldChange }\SpecialCharTok{\textgreater{}} \DecValTok{1} \SpecialCharTok{|}\NormalTok{ WT\_Case\_DKO\_Case}\SpecialCharTok{$}\NormalTok{log2FoldChange }\SpecialCharTok{\textless{}} \SpecialCharTok{{-}}\DecValTok{1}\NormalTok{), ]), }\StringTok{" DE genes in WT vs DKO case comparison having padj \textless{} 0.05 and | log2FC | \textgreater{} 1."}\NormalTok{)}
\end{Highlighting}
\end{Shaded}

\begin{verbatim}
## [1] "There has been found 3972 DE genes in WT vs DKO case comparison having padj < 0.05 and | log2FC | > 1."
\end{verbatim}

\begin{Shaded}
\begin{Highlighting}[]
\FunctionTok{paste0}\NormalTok{(}\StringTok{"There has been found "}\NormalTok{, }\FunctionTok{nrow}\NormalTok{(WT\_Control\_WT\_Case), }\StringTok{" DE genes in WT control vs case comparison having padj \textless{} 0.05."}\NormalTok{)}
\end{Highlighting}
\end{Shaded}

\begin{verbatim}
## [1] "There has been found 11055 DE genes in WT control vs case comparison having padj < 0.05."
\end{verbatim}

\begin{Shaded}
\begin{Highlighting}[]
\FunctionTok{paste0}\NormalTok{(}\StringTok{"There has been found "}\NormalTok{, }\FunctionTok{nrow}\NormalTok{(WT\_Control\_WT\_Case[}\FunctionTok{which}\NormalTok{(WT\_Control\_WT\_Case}\SpecialCharTok{$}\NormalTok{log2FoldChange }\SpecialCharTok{\textgreater{}} \DecValTok{1} \SpecialCharTok{|}\NormalTok{ WT\_Control\_WT\_Case}\SpecialCharTok{$}\NormalTok{log2FoldChange }\SpecialCharTok{\textless{}} \SpecialCharTok{{-}}\DecValTok{1}\NormalTok{), ]), }\StringTok{" DE genes in WT control vs case comparison having padj \textless{} 0.05 and | log2FC | \textgreater{} 1."}\NormalTok{)}
\end{Highlighting}
\end{Shaded}

\begin{verbatim}
## [1] "There has been found 6766 DE genes in WT control vs case comparison having padj < 0.05 and | log2FC | > 1."
\end{verbatim}

\begin{Shaded}
\begin{Highlighting}[]
\FunctionTok{paste0}\NormalTok{(}\StringTok{"There has been found "}\NormalTok{, }\FunctionTok{nrow}\NormalTok{(DKO\_Control\_DKO\_Case), }\StringTok{" DE genes in DKO control vs case comparison having padj \textless{} 0.05."}\NormalTok{)}
\end{Highlighting}
\end{Shaded}

\begin{verbatim}
## [1] "There has been found 11405 DE genes in DKO control vs case comparison having padj < 0.05."
\end{verbatim}

\begin{Shaded}
\begin{Highlighting}[]
\FunctionTok{paste0}\NormalTok{(}\StringTok{"There has been found "}\NormalTok{, }\FunctionTok{nrow}\NormalTok{(DKO\_Control\_DKO\_Case[}\FunctionTok{which}\NormalTok{(DKO\_Control\_DKO\_Case}\SpecialCharTok{$}\NormalTok{log2FoldChange }\SpecialCharTok{\textgreater{}} \DecValTok{1} \SpecialCharTok{|}\NormalTok{ DKO\_Control\_DKO\_Case}\SpecialCharTok{$}\NormalTok{log2FoldChange }\SpecialCharTok{\textless{}} \SpecialCharTok{{-}}\DecValTok{1}\NormalTok{), ]), }\StringTok{" DE genes in DKO control vs case comparison having padj \textless{} 0.05 and | log2FC | \textgreater{} 1."}\NormalTok{)}
\end{Highlighting}
\end{Shaded}

\begin{verbatim}
## [1] "There has been found 6785 DE genes in DKO control vs case comparison having padj < 0.05 and | log2FC | > 1."
\end{verbatim}

How many of the DE genes are up-regulated vs down-regulated meaningfully
(padj \textless{} 0.05)?

\begin{Shaded}
\begin{Highlighting}[]
\FunctionTok{paste0}\NormalTok{(}\StringTok{"For WT vs DKO case are "}\NormalTok{, }\FunctionTok{sum}\NormalTok{(WT\_Case\_DKO\_Case}\SpecialCharTok{$}\NormalTok{log2FoldChange }\SpecialCharTok{\textgreater{}} \DecValTok{1}\NormalTok{), }\StringTok{" upregulated and "}\NormalTok{, }\FunctionTok{length}\NormalTok{(WT\_Case\_DKO\_Case}\SpecialCharTok{$}\NormalTok{log2FoldChange) }\SpecialCharTok{{-}} \FunctionTok{sum}\NormalTok{(WT\_Case\_DKO\_Case}\SpecialCharTok{$}\NormalTok{log2FoldChange }\SpecialCharTok{\textgreater{}} \SpecialCharTok{{-}}\DecValTok{1}\NormalTok{), }\StringTok{" downregulated."}\NormalTok{)}
\end{Highlighting}
\end{Shaded}

\begin{verbatim}
## [1] "For WT vs DKO case are 1909 upregulated and 2063 downregulated."
\end{verbatim}

\begin{Shaded}
\begin{Highlighting}[]
\FunctionTok{paste0}\NormalTok{(}\StringTok{"For WT controls vs cases are "}\NormalTok{, }\FunctionTok{sum}\NormalTok{(WT\_Control\_WT\_Case}\SpecialCharTok{$}\NormalTok{log2FoldChange }\SpecialCharTok{\textgreater{}} \DecValTok{1}\NormalTok{), }\StringTok{" upregulated and "}\NormalTok{, }\FunctionTok{length}\NormalTok{(WT\_Control\_WT\_Case}\SpecialCharTok{$}\NormalTok{log2FoldChange) }\SpecialCharTok{{-}} \FunctionTok{sum}\NormalTok{(WT\_Control\_WT\_Case}\SpecialCharTok{$}\NormalTok{log2FoldChange }\SpecialCharTok{\textgreater{}} \SpecialCharTok{{-}}\DecValTok{1}\NormalTok{), }\StringTok{" downregulated."}\NormalTok{)}
\end{Highlighting}
\end{Shaded}

\begin{verbatim}
## [1] "For WT controls vs cases are 3961 upregulated and 2805 downregulated."
\end{verbatim}

\begin{Shaded}
\begin{Highlighting}[]
\FunctionTok{paste0}\NormalTok{(}\StringTok{"For DKO control vs case are "}\NormalTok{, }\FunctionTok{sum}\NormalTok{(DKO\_Control\_DKO\_Case}\SpecialCharTok{$}\NormalTok{log2FoldChange }\SpecialCharTok{\textgreater{}} \DecValTok{1}\NormalTok{), }\StringTok{" upregulated and "}\NormalTok{, }\FunctionTok{length}\NormalTok{(DKO\_Control\_DKO\_Case}\SpecialCharTok{$}\NormalTok{log2FoldChange) }\SpecialCharTok{{-}} \FunctionTok{sum}\NormalTok{(DKO\_Control\_DKO\_Case}\SpecialCharTok{$}\NormalTok{log2FoldChange }\SpecialCharTok{\textgreater{}} \SpecialCharTok{{-}}\DecValTok{1}\NormalTok{), }\StringTok{" downregulated."}\NormalTok{)}
\end{Highlighting}
\end{Shaded}

\begin{verbatim}
## [1] "For DKO control vs case are 3874 upregulated and 2911 downregulated."
\end{verbatim}

Based on the original publication, select 2-3 genes that are of
particular interest and investigate their expression level. make
\textbf{SUPPLEMENTARY REPORT TABLE 2} for the selected genes with name,
ENSEMBL name, pathway, count, where DE.

\begin{verbatim}
common <- rownames(WT_typeII)[which(! rownames(WT_typeII) %in% rownames(DKO_typeII))]

DKO_typeI[which(rownames(DKO_typeI) == "ENSMUSG00000052776"),]
\end{verbatim}

After comparing the DE genes in the three different comparisons, I
choose four representative ones from different BP families that explain
best what has been concluded in the paper.

\begin{itemize}
\tightlist
\item
  Oas1a (type I) ENSMUSG00000052776 present in all three comparisons, so
  this gene function is related to the disease and not influenced by the
  DKO
\item
  Ifit1 (type I) ENSMUSG00000034459 not present in WT control and case,
  but not present in DKO case, so affected by DKO
\item
  Gbp5 (type II) ENSMUSG00000105504 \textbar\textbar{}
\item
  Irf1 (type II) ENSMUSG00000018899 \textbar\textbar{}
\end{itemize}

Oas1a (type I) ENSMUSG00000052776 Ifit1 (type I) ENSMUSG00000034459 Gbp5
(type II) ENSMUSG00000105504 Irf1 (type II) ENSMUSG00000018899

take the vsd information on the corrected variance and plot those

\begin{Shaded}
\begin{Highlighting}[]
\CommentTok{\# Define genes of interest}
\NormalTok{gene\_annot }\OtherTok{\textless{}{-}} \FunctionTok{data.frame}\NormalTok{(}
  \AttributeTok{Gene =} \FunctionTok{c}\NormalTok{(}
    \StringTok{"ENSMUSG00000052776"}\NormalTok{,}
    \StringTok{"ENSMUSG00000034459"}\NormalTok{,}
    \StringTok{"ENSMUSG00000105504"}\NormalTok{,}
    \StringTok{"ENSMUSG00000018899"}
\NormalTok{  ),}
  \AttributeTok{Label =} \FunctionTok{c}\NormalTok{(}
    \StringTok{"Oas1a (type I)}\SpecialCharTok{\textbackslash{}n}\StringTok{ENSMUSG00000052776"}\NormalTok{,}
    \StringTok{"Ifit1 (type I)}\SpecialCharTok{\textbackslash{}n}\StringTok{ENSMUSG00000034459"}\NormalTok{,}
    \StringTok{"Gbp5 (type II)}\SpecialCharTok{\textbackslash{}n}\StringTok{ENSMUSG00000105504"}\NormalTok{,}
    \StringTok{"Irf1 (type II)}\SpecialCharTok{\textbackslash{}n}\StringTok{ENSMUSG00000018899"}
\NormalTok{  )}
\NormalTok{)}

\CommentTok{\# Extract expression data for selected genes from vsd\_mat}
\NormalTok{genes\_of\_interest }\OtherTok{\textless{}{-}}\NormalTok{ gene\_annot}\SpecialCharTok{$}\NormalTok{Gene}
\NormalTok{vsd\_subset }\OtherTok{\textless{}{-}}\NormalTok{ vsd\_mat[genes\_of\_interest, ]}

\CommentTok{\# Convert to data frame and add Gene column}
\NormalTok{df\_wide }\OtherTok{\textless{}{-}} \FunctionTok{as.data.frame}\NormalTok{(}\FunctionTok{t}\NormalTok{(vsd\_subset))}
\NormalTok{df\_wide}\SpecialCharTok{$}\NormalTok{Sample }\OtherTok{\textless{}{-}} \FunctionTok{rownames}\NormalTok{(df\_wide)}

\CommentTok{\# Reshape to long format}
\FunctionTok{library}\NormalTok{(tidyr)}
\FunctionTok{library}\NormalTok{(dplyr)}
\end{Highlighting}
\end{Shaded}

\begin{verbatim}
## 
## Attaching package: 'dplyr'
\end{verbatim}

\begin{verbatim}
## The following objects are masked from 'package:stats':
## 
##     filter, lag
\end{verbatim}

\begin{verbatim}
## The following objects are masked from 'package:base':
## 
##     intersect, setdiff, setequal, union
\end{verbatim}

\begin{Shaded}
\begin{Highlighting}[]
\NormalTok{df\_long }\OtherTok{\textless{}{-}}\NormalTok{ df\_wide }\SpecialCharTok{\%\textgreater{}\%}
  \FunctionTok{pivot\_longer}\NormalTok{(}\AttributeTok{cols =} \SpecialCharTok{{-}}\NormalTok{Sample, }\AttributeTok{names\_to =} \StringTok{"Gene"}\NormalTok{, }\AttributeTok{values\_to =} \StringTok{"Expression"}\NormalTok{) }\SpecialCharTok{\%\textgreater{}\%}
  \FunctionTok{left\_join}\NormalTok{(sample\_info, }\AttributeTok{by =} \StringTok{"Sample"}\NormalTok{) }\SpecialCharTok{\%\textgreater{}\%}
  \FunctionTok{left\_join}\NormalTok{(gene\_annot, }\AttributeTok{by =} \StringTok{"Gene"}\NormalTok{)}

\CommentTok{\# Convert Label to a factor with the specific order}
\NormalTok{df\_long}\SpecialCharTok{$}\NormalTok{Label }\OtherTok{\textless{}{-}} \FunctionTok{factor}\NormalTok{(df\_long}\SpecialCharTok{$}\NormalTok{Label, }\AttributeTok{levels =} \FunctionTok{c}\NormalTok{(}
  \StringTok{"Oas1a (type I)}\SpecialCharTok{\textbackslash{}n}\StringTok{ENSMUSG00000052776"}\NormalTok{,}
  \StringTok{"Ifit1 (type I)}\SpecialCharTok{\textbackslash{}n}\StringTok{ENSMUSG00000034459"}\NormalTok{,}
  \StringTok{"Gbp5 (type II)}\SpecialCharTok{\textbackslash{}n}\StringTok{ENSMUSG00000105504"}\NormalTok{,}
  \StringTok{"Irf1 (type II)}\SpecialCharTok{\textbackslash{}n}\StringTok{ENSMUSG00000018899"}
\NormalTok{))}
\end{Highlighting}
\end{Shaded}

\begin{Shaded}
\begin{Highlighting}[]
\CommentTok{\# Define comparisons for statistical testing}
\NormalTok{comparisons }\OtherTok{\textless{}{-}} \FunctionTok{list}\NormalTok{(}
  \FunctionTok{c}\NormalTok{(}\StringTok{"Lung\_WT\_Control"}\NormalTok{, }\StringTok{"Lung\_WT\_Case"}\NormalTok{),}
  \FunctionTok{c}\NormalTok{(}\StringTok{"Lung\_DKO\_Control"}\NormalTok{, }\StringTok{"Lung\_DKO\_Case"}\NormalTok{),}
  \FunctionTok{c}\NormalTok{(}\StringTok{"Lung\_WT\_Case"}\NormalTok{, }\StringTok{"Lung\_DKO\_Case"}\NormalTok{)}
\NormalTok{)}

\FunctionTok{library}\NormalTok{(ggplot2)}
\end{Highlighting}
\end{Shaded}

\begin{verbatim}
## Warning: package 'ggplot2' was built under R version 4.5.2
\end{verbatim}

\begin{Shaded}
\begin{Highlighting}[]
\FunctionTok{library}\NormalTok{(ggpubr)}

\CommentTok{\# These colors help distinguish between WT and DKO groups clearly [1, 2]}
\NormalTok{group\_colors }\OtherTok{\textless{}{-}} \FunctionTok{c}\NormalTok{(}
  \StringTok{"Lung\_WT\_Control"}  \OtherTok{=} \StringTok{"\#CCCCCC"}\NormalTok{, }\CommentTok{\# Light Grey}
  \StringTok{"Lung\_WT\_Case"}     \OtherTok{=} \StringTok{"\#00008B"}\NormalTok{, }\CommentTok{\# Dark Blue}
  \StringTok{"Lung\_DKO\_Control"} \OtherTok{=} \StringTok{"\#FFD700"}\NormalTok{, }\CommentTok{\# Gold/Yellow}
  \StringTok{"Lung\_DKO\_Case"}    \OtherTok{=} \StringTok{"\#FF4500"}  \CommentTok{\# Orange{-}Red}
\NormalTok{)}

\CommentTok{\# 2. Updated Plotting Code}
\FunctionTok{ggplot}\NormalTok{(df\_long, }\FunctionTok{aes}\NormalTok{(}\AttributeTok{x =}\NormalTok{ Group, }\AttributeTok{y =}\NormalTok{ Expression, }\AttributeTok{fill =}\NormalTok{ Group)) }\SpecialCharTok{+}
  \CommentTok{\# Add transparency (alpha) to the boxplot to make it look cleaner}
  \FunctionTok{geom\_boxplot}\NormalTok{(}\AttributeTok{outlier.shape =} \ConstantTok{NA}\NormalTok{, }\AttributeTok{alpha =} \FloatTok{0.7}\NormalTok{, }\AttributeTok{color =} \StringTok{"black"}\NormalTok{) }\SpecialCharTok{+} 
  \CommentTok{\# Increase jitter transparency to avoid clutter while showing every replicate [5]}
  \FunctionTok{geom\_jitter}\NormalTok{(}\AttributeTok{width =} \FloatTok{0.15}\NormalTok{, }\AttributeTok{size =} \FloatTok{1.5}\NormalTok{, }\AttributeTok{alpha =} \FloatTok{0.5}\NormalTok{) }\SpecialCharTok{+}

  \CommentTok{\# Maintain your factor{-}based ordering for Oas1a, Ifit1, Gbp5, and Irf1 [6, 7]}
  \FunctionTok{facet\_wrap}\NormalTok{(}\SpecialCharTok{\textasciitilde{}}\NormalTok{ Label, }\AttributeTok{scales =} \StringTok{"free\_y"}\NormalTok{) }\SpecialCharTok{+}

  \FunctionTok{stat\_compare\_means}\NormalTok{(}
    \AttributeTok{comparisons =}\NormalTok{ comparisons,}
    \AttributeTok{method =} \StringTok{"wilcox.test"}\NormalTok{,}
    \AttributeTok{label =} \StringTok{"p.signif"}\NormalTok{,}
    \AttributeTok{hide.ns =} \ConstantTok{TRUE}\NormalTok{,}
    \AttributeTok{step.increase =} \FloatTok{0.12} \CommentTok{\# Prevents p{-}value labels from overlapping}
\NormalTok{  ) }\SpecialCharTok{+}

  \CommentTok{\# Switch to theme\_minimal to remove the heavy grey background [3]}
  \FunctionTok{theme\_minimal}\NormalTok{(}\AttributeTok{base\_size =} \DecValTok{12}\NormalTok{) }\SpecialCharTok{+} 
  \FunctionTok{theme}\NormalTok{(}
    \CommentTok{\# Remove the grey boxes behind the facet labels (gene names)}
    \AttributeTok{strip.background =} \FunctionTok{element\_blank}\NormalTok{(), }
    \AttributeTok{strip.text =} \FunctionTok{element\_text}\NormalTok{(}\AttributeTok{face =} \StringTok{"bold"}\NormalTok{, }\AttributeTok{size =} \DecValTok{10}\NormalTok{),}
    
    \CommentTok{\# Clean up the axes and remove background clutter}
    \AttributeTok{panel.grid.minor =} \FunctionTok{element\_blank}\NormalTok{(),}
    \AttributeTok{panel.grid.major.x =} \FunctionTok{element\_blank}\NormalTok{(),}
    \AttributeTok{axis.line =} \FunctionTok{element\_line}\NormalTok{(}\AttributeTok{color =} \StringTok{"black"}\NormalTok{),}
    \AttributeTok{axis.text.x =} \FunctionTok{element\_text}\NormalTok{(}\AttributeTok{angle =} \DecValTok{45}\NormalTok{, }\AttributeTok{hjust =} \DecValTok{1}\NormalTok{),}
    
    \AttributeTok{legend.position =} \StringTok{"none"} \CommentTok{\# Groups are already labeled on the X{-}axis}
\NormalTok{  ) }\SpecialCharTok{+}
  
  \CommentTok{\# Apply the custom color scale}
  \FunctionTok{scale\_fill\_manual}\NormalTok{(}\AttributeTok{values =}\NormalTok{ group\_colors) }\SpecialCharTok{+}
  
  \FunctionTok{labs}\NormalTok{(}
    \AttributeTok{title =} \StringTok{"VST Expression of Selected Interferon{-}Related Genes"}\NormalTok{,}
    \AttributeTok{y =} \StringTok{"VST Normalized Expression"}\NormalTok{,}
    \AttributeTok{x =} \ConstantTok{NULL}
\NormalTok{  )}
\end{Highlighting}
\end{Shaded}

\pandocbounded{\includegraphics[keepaspectratio]{DEshort_files/figure-latex/unnamed-chunk-14-1.pdf}}

\begin{Shaded}
\begin{Highlighting}[]
\FunctionTok{ggsave}\NormalTok{(}
  \AttributeTok{filename =} \StringTok{"IFN\_genes\_VST\_boxplot.png"}\NormalTok{, }\CommentTok{\# You can also use .pdf, .jpeg, or .tiff}
  \AttributeTok{width =} \DecValTok{12}\NormalTok{,                          }\CommentTok{\# Width in inches}
  \AttributeTok{height =} \FloatTok{8.5}\NormalTok{,                          }\CommentTok{\# Height in inches}
  \AttributeTok{dpi =} \DecValTok{300}                            \CommentTok{\# High resolution for the report}
\NormalTok{)}
\end{Highlighting}
\end{Shaded}

\textbf{genes from paper}

\begin{Shaded}
\begin{Highlighting}[]
\CommentTok{\# Install database}
\ControlFlowTok{if}\NormalTok{ (}\SpecialCharTok{!}\FunctionTok{require}\NormalTok{(}\StringTok{"org.Mm.eg.db"}\NormalTok{, }\AttributeTok{quietly =} \ConstantTok{TRUE}\NormalTok{))}
\NormalTok{  BiocManager}\SpecialCharTok{::}\FunctionTok{install}\NormalTok{(}\StringTok{"org.Mm.eg.db"}\NormalTok{)}

\FunctionTok{library}\NormalTok{(org.Mm.eg.db)}
\FunctionTok{library}\NormalTok{(AnnotationDbi)}

\CommentTok{\# function to obtain ENSEMBL names from their name}
\NormalTok{gfilter }\OtherTok{\textless{}{-}} \ControlFlowTok{function}\NormalTok{(symbols, dss)\{}
\NormalTok{  ens\_ids }\OtherTok{\textless{}{-}} \FunctionTok{mapIds}\NormalTok{(}
\NormalTok{    org.Mm.eg.db,}
    \AttributeTok{keys      =}\NormalTok{ symbols,}
    \AttributeTok{column    =} \StringTok{"ENSEMBL"}\NormalTok{,}
    \AttributeTok{keytype   =} \StringTok{"SYMBOL"}\NormalTok{,}
    \AttributeTok{multiVals =} \StringTok{"first"}
\NormalTok{  )}
  
\NormalTok{  dss\_fil }\OtherTok{\textless{}{-}}\NormalTok{ dss[}\FunctionTok{grepl}\NormalTok{(}\FunctionTok{paste}\NormalTok{(}\FunctionTok{unname}\NormalTok{(ens\_ids), }\AttributeTok{collapse=}\StringTok{"|"}\NormalTok{), }\FunctionTok{row.names}\NormalTok{(dss), }\AttributeTok{ignore.case =} \ConstantTok{TRUE}\NormalTok{), ]}
  
  \CommentTok{\# add gene name to the data frame}
\NormalTok{  gene\_map\_one }\OtherTok{\textless{}{-}} \FunctionTok{data.frame}\NormalTok{(}
    \AttributeTok{ENSEMBL =} \FunctionTok{unname}\NormalTok{(ens\_ids),}
    \AttributeTok{name    =} \FunctionTok{names}\NormalTok{(ens\_ids),}
    \AttributeTok{stringsAsFactors =} \ConstantTok{FALSE}
\NormalTok{  )}

\NormalTok{  dss\_fil}\SpecialCharTok{$}\NormalTok{name }\OtherTok{\textless{}{-}}\NormalTok{ gene\_map\_one}\SpecialCharTok{$}\NormalTok{name[}\FunctionTok{match}\NormalTok{(}\FunctionTok{rownames}\NormalTok{(dss\_fil), gene\_map\_one}\SpecialCharTok{$}\NormalTok{ENSEMBL)]}
  
  \FunctionTok{return}\NormalTok{(dss\_fil)}
\NormalTok{\}}

\CommentTok{\# list of gene names to convert into ENSEMBLE codes}
\DocumentationTok{\#\# L5 (type I IFN)}
\NormalTok{symbols }\OtherTok{\textless{}{-}} \FunctionTok{c}\NormalTok{(}
  \StringTok{"Ifnb1"}\NormalTok{, }\StringTok{"Ifit1"}\NormalTok{, }\StringTok{"Ifit3"}\NormalTok{, }\StringTok{"Oas1a"}\NormalTok{, }\StringTok{"Oas2"}\NormalTok{, }\StringTok{"Oas3"}\NormalTok{, }\StringTok{"Oasl1"}\NormalTok{, }\StringTok{"Oasl2"}\NormalTok{, }\StringTok{"Mx1"}\NormalTok{, }\StringTok{"Stat2"}\NormalTok{, }\StringTok{"Irf7"}\NormalTok{, }\StringTok{"Irf9"}
\NormalTok{)}

\NormalTok{Case\_typeI }\OtherTok{\textless{}{-}} \FunctionTok{gfilter}\NormalTok{(symbols, WT\_Case\_DKO\_Case)}
\NormalTok{WT\_typeI }\OtherTok{\textless{}{-}} \FunctionTok{gfilter}\NormalTok{(symbols, WT\_Control\_WT\_Case)}
\NormalTok{DKO\_typeI }\OtherTok{\textless{}{-}} \FunctionTok{gfilter}\NormalTok{(symbols, DKO\_Control\_DKO\_Case)}


\DocumentationTok{\#\# L7 (type II IFN)}
\NormalTok{symbols }\OtherTok{\textless{}{-}} \FunctionTok{c}\NormalTok{(}
  \StringTok{"Ifng"}\NormalTok{, }\StringTok{"Irf1"}\NormalTok{, }\StringTok{"Irf2"}\NormalTok{, }\StringTok{"Irf8"}\NormalTok{, }\StringTok{"Il12rb1"}\NormalTok{, }\StringTok{"Il12rb2"}\NormalTok{, }\StringTok{"Tap1"}\NormalTok{, }\StringTok{"Tap2"}\NormalTok{, }\StringTok{"Gbp2"}\NormalTok{, }\StringTok{"Gbp3"}\NormalTok{, }\StringTok{"Gbp4"}\NormalTok{, }\StringTok{"Gbp5"}\NormalTok{, }\StringTok{"Gbp6"}\NormalTok{, }\StringTok{"Gbp8"}\NormalTok{, }\StringTok{"Gbp9"}\NormalTok{, }\StringTok{"Gbp10"}\NormalTok{, }\StringTok{"Gbp11"}
\NormalTok{)}

\NormalTok{Case\_typeII }\OtherTok{\textless{}{-}} \FunctionTok{gfilter}\NormalTok{(symbols, WT\_Case\_DKO\_Case)}
\NormalTok{WT\_typeII }\OtherTok{\textless{}{-}} \FunctionTok{gfilter}\NormalTok{(symbols, WT\_Control\_WT\_Case)}
\NormalTok{DKO\_typeII }\OtherTok{\textless{}{-}} \FunctionTok{gfilter}\NormalTok{(symbols, DKO\_Control\_DKO\_Case)}
\end{Highlighting}
\end{Shaded}

\begin{center}\rule{0.5\linewidth}{0.5pt}\end{center}

\subsection{7. Overrepresentation
analysis}\label{overrepresentation-analysis}

\textbf{Hypergeometric test} is used to check overrepresentation for
each biological pathways (BP).

Install clusterProfiler

\begin{Shaded}
\begin{Highlighting}[]
\ControlFlowTok{if}\NormalTok{ (}\SpecialCharTok{!}\FunctionTok{require}\NormalTok{(}\StringTok{"BiocManager"}\NormalTok{, }\AttributeTok{quietly =} \ConstantTok{TRUE}\NormalTok{))}
  \FunctionTok{install.packages}\NormalTok{(}\StringTok{"BiocManager"}\NormalTok{)}

\ControlFlowTok{if}\NormalTok{ (}\SpecialCharTok{!}\FunctionTok{require}\NormalTok{(}\StringTok{"clusterProfiler"}\NormalTok{, }\AttributeTok{quietly =} \ConstantTok{TRUE}\NormalTok{))}
\NormalTok{  BiocManager}\SpecialCharTok{::}\FunctionTok{install}\NormalTok{(}\StringTok{"clusterProfiler"}\NormalTok{)}
\CommentTok{\# the link to enrichplot is broken in bioconductor, install directly }
\ControlFlowTok{if}\NormalTok{ (}\SpecialCharTok{!}\FunctionTok{require}\NormalTok{(}\StringTok{"enrichplot"}\NormalTok{, }\AttributeTok{quietly =} \ConstantTok{TRUE}\NormalTok{))}
\NormalTok{  BiocManager}\SpecialCharTok{::}\FunctionTok{install}\NormalTok{(}\StringTok{"enrichplot"}\NormalTok{, }\AttributeTok{force =} \ConstantTok{TRUE}\NormalTok{) }
\CommentTok{\# without updating the other packages (might create problems with dependencies)}

\FunctionTok{library}\NormalTok{(}\StringTok{"clusterProfiler"}\NormalTok{)}
\end{Highlighting}
\end{Shaded}

Download Bioconductor package with genome wide annotation for your
species

\begin{Shaded}
\begin{Highlighting}[]
\ControlFlowTok{if}\NormalTok{ (}\SpecialCharTok{!}\FunctionTok{require}\NormalTok{(}\StringTok{"org.Mm.eg.db"}\NormalTok{, }\AttributeTok{quietly =} \ConstantTok{TRUE}\NormalTok{))}
\NormalTok{  BiocManager}\SpecialCharTok{::}\FunctionTok{install}\NormalTok{(}\StringTok{"org.Mm.eg.db"}\NormalTok{)}

\FunctionTok{library}\NormalTok{(}\StringTok{"org.Mm.eg.db"}\NormalTok{)}
\end{Highlighting}
\end{Shaded}

Call enrichGO

\begin{Shaded}
\begin{Highlighting}[]
\CommentTok{\# Select significant genes WT control and cases and run the enrichment via Gene Ontology}
\NormalTok{ego\_gen }\OtherTok{\textless{}{-}} \ControlFlowTok{function}\NormalTok{(comparison, DEresults)\{}
\NormalTok{  ego }\OtherTok{\textless{}{-}}\NormalTok{ clusterProfiler}\SpecialCharTok{::}\FunctionTok{enrichGO}\NormalTok{(}
    \AttributeTok{gene =} \FunctionTok{rownames}\NormalTok{(comparison), }
    \AttributeTok{universe =} \FunctionTok{rownames}\NormalTok{(DEresults),}
    \AttributeTok{OrgDb =}\NormalTok{ org.Mm.eg.db,}
    \AttributeTok{ont =} \StringTok{"BP"}\NormalTok{,}
    \AttributeTok{keyType =} \StringTok{"ENSEMBL"}\NormalTok{,}
    \AttributeTok{maxGSSize =} \DecValTok{150}\NormalTok{,}
    \AttributeTok{pvalueCutoff =} \FloatTok{0.05}
\NormalTok{    )}
  
  \CommentTok{\# filter for genes with log2FC higher than +/{-} 1 (2x DE)}
\NormalTok{  ego}\SpecialCharTok{@}\NormalTok{result }\OtherTok{\textless{}{-}}\NormalTok{ ego}\SpecialCharTok{@}\NormalTok{result[}\FunctionTok{which}\NormalTok{(}\SpecialCharTok{{-}}\DecValTok{1} \SpecialCharTok{\textgreater{}}\NormalTok{ ego}\SpecialCharTok{@}\NormalTok{result}\SpecialCharTok{$}\NormalTok{FoldEnrichment }\SpecialCharTok{|} \DecValTok{1} \SpecialCharTok{\textless{}}\NormalTok{ ego}\SpecialCharTok{@}\NormalTok{result}\SpecialCharTok{$}\NormalTok{FoldEnrichment), ]}
  
  \FunctionTok{return}\NormalTok{(ego)}
\NormalTok{\}}

\NormalTok{ego\_WT }\OtherTok{\textless{}{-}} \FunctionTok{ego\_gen}\NormalTok{(WT\_Control\_WT\_Case, dss)}
\NormalTok{ego\_DKO }\OtherTok{\textless{}{-}} \FunctionTok{ego\_gen}\NormalTok{(DKO\_Control\_DKO\_Case, dss)}
\NormalTok{ego\_Case }\OtherTok{\textless{}{-}} \FunctionTok{ego\_gen}\NormalTok{(WT\_Case\_DKO\_Case, dss)}

\CommentTok{\# plotting the results}
\FunctionTok{dotplot}\NormalTok{(ego\_Case)}
\end{Highlighting}
\end{Shaded}

\pandocbounded{\includegraphics[keepaspectratio]{DEshort_files/figure-latex/unnamed-chunk-18-1.pdf}}

\begin{Shaded}
\begin{Highlighting}[]
\FunctionTok{dotplot}\NormalTok{(ego\_DKO)}
\end{Highlighting}
\end{Shaded}

\pandocbounded{\includegraphics[keepaspectratio]{DEshort_files/figure-latex/unnamed-chunk-18-2.pdf}}

\begin{Shaded}
\begin{Highlighting}[]
\FunctionTok{dotplot}\NormalTok{(ego\_WT)}
\end{Highlighting}
\end{Shaded}

\pandocbounded{\includegraphics[keepaspectratio]{DEshort_files/figure-latex/unnamed-chunk-18-3.pdf}}

\begin{itemize}
\item
  \emph{gene = sig}

  sig should be a vector of ENSEMBL gene IDs that are significantly DE.
\item
  \emph{universe = rownames(mat)}

  All genes measured in the experiment. Since mat is your DESeq2 input
  matrix, its row names are all ENSEMBL IDs → perfect universe.
\item
  \emph{OrgDb = org.Mm.eg.db}

  Genome-wide mouse annotation package from Bioconductor.
\item
  \emph{ont = ``BP''} or ``MF'', ``CC'', or ``ALL''

  GO subontology:

  \begin{itemize}
  \tightlist
  \item
    BP: Biological Process
  \item
    MF: Molecular Function
  \item
    CC: Cellular Component
  \end{itemize}
\item
  \emph{keyType = ``ENSEMBL''}

  Because the gene IDs (row names) are ENSEMBL IDs like
  ``ENSMUSG00000000001''.
\item
  \emph{maxGSSize = 150} maximal size of genes annotated for testing
\end{itemize}

\end{document}
